% Trigonometrie · Teildreiecke und Vermessung · Prüfungs-Fokus MSA – bank/trigonometrie/e3.jsonl (zusammenbau v0.6, Prüfungs-Fokus)
\einheitenkopf[e3]{Einheit 3 · Rechtwinklige Teildreiecke in Figuren und Vermessung}
\zweigzeile{ab Kl. 10 (am Gymnasium ab Kl. 9) · P10 ×12}
\setcounter{aufgabe}{0}

% Anlauf (ohne Prüfkennung)
\begin{aufgabe}{Teildreiecke und Vermessung – Anlauf}
\begin{teile}
\teil Fahre im Trapez das rechtwinklige Teildreieck mit der Höhe $h$ links nach. Beschrifte seine Seiten mit H, G, A – vom Basiswinkel unten links aus.

\trapez[hoehe=h]{5}{3}{2.5}
\teil Gleichschenkliges Dreieck: Schenkel $7$ cm, Basiswinkel $63^\circ$. Wie lang ist die Höhe $h$ auf die Grundseite? h ≈ \leerfeld[cm]
\teil Trapez: Höhe $4$ m, Basiswinkel $62^\circ$. Wie lang ist der Schenkel $s$? s ≈ \leerfeld[m]
\end{teile}
\end{aufgabe}

\begin{aufgabe}{Teildreiecke und Vermessung – Prüfungsaufgaben}
\begin{teile}
\teil Ein Messgerät steht $1{,}4$ m über dem Boden. Die Sichtlinie zur Hügelkuppe ist $96{,}5$ m lang und steigt unter $27^\circ$ an. Wie hoch ist der Hügel? Auf der Kuppe steht ein $9$ m hoher Aussichtsplattform – wie hoch über dem Boden liegt seine Oberkante? \hfill (P18) \\ \leerfeld[m] ; \leerfeld[m]
\teil Ein Messgerät steht $1{,}7$ m über dem Boden. Die Sichtlinie zur Hügelkuppe ist $148$ m lang und steigt unter $24^\circ$ an. Wie hoch ist der Hügel? Auf der Kuppe steht ein $7{,}5$ m hoher Funkmast – wie hoch über dem Boden liegt seine Oberkante? \hfill (P18) \\ \leerfeld[m] ; \leerfeld[m]
\teil Ein rechtwinkliges Trapez hat die senkrechten Seiten $18$ cm und $7$ cm und die waagerechte Grundseite $16$ cm. $\alpha$ ist der Winkel an der oberen Ecke der $18$-cm-Seite zwischen dieser Seite und der schrägen Seite. Berechne $\alpha$. \hfill (P20) \\ $\alpha \approx$ \leerfeld $^\circ$
\teil Ein rechtwinkliges Trapez hat die senkrechten Seiten $13$ cm und $5$ cm und die waagerechte Grundseite $10$ cm. $\alpha$ ist der Winkel an der oberen Ecke der $13$-cm-Seite zwischen dieser Seite und der schrägen Seite. Berechne $\alpha$. \hfill (P20) \\ $\alpha \approx$ \leerfeld $^\circ$
\teil Im Dreieck $ABC$ liegt $F$ auf $AC$, $BF$ steht senkrecht auf $AC$. $AB = 96$ cm, der Winkel bei $A$ ist $41^\circ$, $\angle FBA = 49^\circ$, $\angle CBF = 57^\circ$. Zeige, dass $BF \approx 63{,}0$ cm und $AF \approx 72{,}5$ cm sind, und berechne $BC$. \hfill (P23) \\ BC ≈ \leerfeld[cm]
\teil Im Dreieck $ABC$ liegt $F$ auf $AC$, $BF$ steht senkrecht auf $AC$. $AB = 86$ cm, der Winkel bei $A$ ist $33^\circ$, $\angle FBA = 57^\circ$, $\angle CBF = 62^\circ$. Zeige, dass $BF \approx 46{,}8$ cm und $AF \approx 72{,}1$ cm sind, und berechne $BC$. \hfill (P23) \\ BC ≈ \leerfeld[cm]
\teil Im Parallelogramm $ABCD$ ist die Höhe $h = CF \approx 21{,}4$ cm bekannt; $F$ liegt auf der Verlängerung von $AB$. Der Winkel $CAB$ beträgt $19^\circ$. Berechne die Länge der Diagonale $AC$. \hfill (P26F) \\ AC ≈ \leerfeld[cm]
\teil Im Parallelogramm $ABCD$ ist die Höhe $h = CF \approx 15{,}6$ cm bekannt; $F$ liegt auf der Verlängerung von $AB$. Der Winkel $CAB$ beträgt $23^\circ$. Berechne die Länge der Diagonale $AC$. \hfill (P26F) \\ AC ≈ \leerfeld[cm]
\steil Im Dreieck $ABC$ ist $AD$ die Höhe auf $BC$ ($D$ zwischen $B$ und $C$), $AD = 610$ m. $\angle ACD = 59^\circ$, $\angle CAB = 72^\circ$. Berechne $BD$. \hfill (P16*) \\ BD ≈ \leerfeld[m]
\steil Im Dreieck $ABC$ ist $AD$ die Höhe auf $BC$ ($D$ zwischen $B$ und $C$), $AD = 380$ m. $\angle ACD = 63^\circ$, $\angle CAB = 71^\circ$. Berechne $BD$. \hfill (P16*) \\ BD ≈ \leerfeld[m]
\teil Im Drachen $ABCD$ schneiden sich die Diagonalen in $E$ rechtwinklig. Bekannt sind $AD$ und der Winkel $\delta$ bei $D$ im Dreieck $AED$. Welche Größen brauchst du, und in welcher Reihenfolge rechnest du, um den Flächeninhalt von $AED$ zu bestimmen? \hfill (P15)
\teil Im Parallelogramm $ABCD$ ist $DF$ die Höhe auf $AB$ ($F$ auf $AB$). Bekannt sind $AD$ und der Winkel $\alpha$ bei $A$. Welche Größen brauchst du, und in welcher Reihenfolge rechnest du, um den Flächeninhalt des Dreiecks $AFD$ zu bestimmen? \hfill (P15)
\end{teile}
\end{aufgabe}

\medskip
\uebersichtskasten{Teildreieck: Suche das rechtwinklige Dreieck in der Figur, fahre es nach und benenne die Seiten vom gegebenen Winkel aus – die Höhe, ein Abstand oder eine Hilfslinie ist meist eine Kathete, die schräge Seite die Hypotenuse. \\ Trapez mit Höhe 5 m und Basiswinkel 50°: Schenkel = 5 : sin 50° ≈ 6,5 m \\ Vermessung: Abstand und Höhenwinkel geben die Höhe über dem Gerät; Gerätehöhe, Sockel oder Aufbau werden danach addiert. \\ Abstand 40 m, Höhenwinkel 35°, Augenhöhe 1,6 m: Turmhöhe = 40 · tan 35° + 1,6 ≈ 28,0 + 1,6 = 29,6 m \\ Zwei Dreiecke: Das Ergebnis des ersten Dreiecks ist die gegebene Seite des zweiten; ein Teilwinkel ist die Differenz zweier Winkel an derselben Ecke.}
